\section{Pendahuluan dan Sejarah Aljabar Boolean}

Aljabar Boolean merupakan cabang matematika yang secara khusus mempelajari struktur aljabar atas himpunan dengan dua elemen, yang umumnya dinotasikan sebagai $\{0, 1\}$ atau $\{False, True\}$. Sistem aljabar ini berbeda secara fundamental dari aljabar konvensional yang beroperasi pada bilangan real, karena aljabar Boolean hanya mengenal dua nilai kebenaran dan menggunakan tiga operasi dasar: penjumlahan logis (OR), perkalian logis (AND), dan komplemen (NOT). Pemahaman mendalam terhadap aljabar Boolean menjadi prasyarat esensial bagi mahasiswa informatika, mengingat seluruh operasi komputasi digital pada dasarnya dibangun di atas fondasi struktur aljabar ini \cite{rosen2019}.

Aljabar Boolean pertama kali diperkenalkan oleh matematikawan Inggris George Boole pada tahun 1854 melalui karya monumentalnya \textit{An Investigation of the Laws of Thought}. Dalam karya tersebut, Boole mengusulkan sebuah sistem aljabar yang mampu merepresentasikan proses penalaran logis secara simbolik dan matematis. Kontribusi Boole tidak hanya bersifat teoretis, melainkan juga membuka jalan bagi formalisasi logika yang sebelumnya hanya dibahas secara filosofis oleh para pemikir seperti Aristoteles dan Leibniz \cite{wiki_boolean}. Meskipun pada masanya gagasan Boole belum menemukan aplikasi praktis yang luas, fondasi yang diletakkannya menjadi cikal bakal seluruh arsitektur komputasi modern.

Titik balik penerapan praktis aljabar Boolean terjadi pada tahun 1938, ketika Claude Elwood Shannon, seorang mahasiswa pascasarjana di Massachusetts Institute of Technology (MIT), mempublikasikan tesis masternya yang berjudul \textit{A Symbolic Analysis of Relay and Switching Circuits} \cite{shannon1938}. Shannon mendemonstrasikan bahwa aljabar Boolean dapat digunakan untuk menganalisis dan merancang rangkaian sakelar (switching circuits) pada sistem telepon. Penemuan ini menunjukkan bahwa setiap fungsi logis dapat direalisasikan secara fisik menggunakan rangkaian relay elektromekanik, dan sebaliknya, setiap rangkaian sakelar dapat dianalisis menggunakan ekspresi aljabar Boolean. Karya Shannon ini secara luas dianggap sebagai salah satu tesis master paling berpengaruh dalam sejarah ilmu pengetahuan, karena menjembatani matematika abstrak dengan rekayasa elektrik.

\begin{catatan}
Kontribusi Shannon memberikan landasan teoretis bagi seluruh desain rangkaian digital modern. Setiap komputer, telepon seluler, dan perangkat elektronik digital yang digunakan saat ini pada dasarnya beroperasi berdasarkan prinsip-prinsip yang pertama kali diformalisasikan oleh Shannon melalui penerapan aljabar Boolean pada rangkaian sakelar \cite{shannon1938}.
\end{catatan}

\subsection{Relasi Isomorfik: Logika, Himpunan, dan Boolean}

Salah satu aspek paling menarik dari aljabar Boolean adalah relasinya yang bersifat \logic{isomorfik} dengan dua struktur matematika lainnya, yaitu logika proposisi dan teori himpunan. Isomorfisme ini berarti bahwa ketiga sistem memiliki struktur aljabar yang identik, sehingga setiap teorema atau hukum yang berlaku di satu sistem dapat diterjemahkan secara langsung ke dua sistem lainnya hanya dengan mengganti simbol operator dan elemen dasarnya \cite{rosen2019}. Pemahaman terhadap isomorfisme ini memberikan wawasan yang mendalam tentang kesatuan fundamental matematika.

Dalam logika proposisi, elemen dasar adalah proposisi yang bernilai \textit{True} (T) atau \textit{False} (F), dengan operator konjungsi ($\wedge$), disjungsi ($\vee$), dan negasi ($\neg$). Dalam teori himpunan, elemen dasar adalah himpunan-himpunan yang beroperasi dengan irisan ($\cap$), gabungan ($\cup$), dan komplemen ($^c$). Sedangkan dalam aljabar Boolean, variabel bernilai 1 atau 0 dengan operator perkalian ($\cdot$), penjumlahan ($+$), dan komplemen ($'$). Ketiga struktur ini saling memetakan satu sama lain sebagaimana diilustrasikan pada Gambar~\ref{fig:isomorfik}.

\begin{figure}[H]
\centering
\begin{tikzpicture}[
    node distance=2.2cm,
    box/.style={rectangle, draw, rounded corners, fill=blue!10, text width=3.8cm, align=center, minimum height=1.6cm},
    arrow/.style={<->, thick, >=stealth}
]
    \node[box] (logic) {\textbf{Logika Proposisi} \\ $p, q$ \\ $\land, \lor, \neg$ \\ Nilai: T, F};
    \node[box, below right=of logic] (set) {\textbf{Teori Himpunan} \\ $A, B$ \\ $\cap, \cup, ^c$ \\ Relasi: $\subseteq$};
    \node[box, below left=of logic] (boolean) {\textbf{Aljabar Boolean} \\ $x, y$ \\ $\cdot, +, '$ \\ Nilai: 1, 0};

    \draw[arrow] (logic) -- node[above right] {Isomorfik} (set);
    \draw[arrow] (logic) -- node[above left] {Isomorfik} (boolean);
    \draw[arrow] (set) -- node[below] {Isomorfik} (boolean);
\end{tikzpicture}
\caption{Hubungan Isomorfik antara Logika Proposisi, Teori Himpunan, dan Aljabar Boolean}
\label{fig:isomorfik}
\end{figure}

Tabel~\ref{tab:korespondensi-tiga-sistem} menyajikan pemetaan lengkap simbol dan konsep antar ketiga sistem tersebut. Dengan tabel ini, mahasiswa dapat dengan mudah mentranslasikan ekspresi dari satu sistem ke sistem lainnya, misalnya hukum De Morgan dalam logika ($\neg(p \wedge q) \equiv \neg p \vee \neg q$) langsung berkorespondensi dengan $(A \cap B)^c = A^c \cup B^c$ pada teori himpunan dan $(x \cdot y)' = x' + y'$ pada aljabar Boolean.

\begin{table}[H]
\centering
\renewcommand{\arraystretch}{1.3}
\begin{tabular}{|l|c|c|c|}
\hline
\textbf{Konsep} & \textbf{Logika Proposisi} & \textbf{Teori Himpunan} & \textbf{Aljabar Boolean} \\ \hline
Nilai benar & True (T) & Semesta ($U$) & 1 \\ \hline
Nilai salah & False (F) & Himpunan kosong ($\emptyset$) & 0 \\ \hline
Operasi AND & Konjungsi ($\wedge$) & Irisan ($\cap$) & Perkalian ($\cdot$) \\ \hline
Operasi OR & Disjungsi ($\vee$) & Gabungan ($\cup$) & Penjumlahan ($+$) \\ \hline
Operasi NOT & Negasi ($\neg$) & Komplemen ($^c$) & Komplemen ($'$) \\ \hline
Relasi subset & Implikasi ($\rightarrow$) & Subset ($\subseteq$) & $x \cdot y' = 0$ \\ \hline
Ekuivalensi & Bikondisional ($\leftrightarrow$) & Kesamaan ($=$) & $x = y$ \\ \hline
\end{tabular}
\caption{Korespondensi Simbol antara Logika Proposisi, Teori Himpunan, dan Aljabar Boolean}
\label{tab:korespondensi-tiga-sistem}
\end{table}

\begin{definisi}[Aljabar Boolean]
Sebuah \logic{aljabar Boolean} adalah sistem matematika $\langle B, +, \cdot, ', 0, 1 \rangle$ yang terdiri atas himpunan $B$ dengan minimal dua elemen berbeda (elemen identitas $0$ dan $1$), dua operator biner --- penjumlahan Boolean ($+$) dan perkalian Boolean ($\cdot$) --- serta satu operator uner komplemen ($'$), yang memenuhi seperangkat aksioma fundamental yang disebut Postulat Huntington \cite{rosen2019}.
\end{definisi}

Kepentingan aljabar Boolean dalam ilmu komputer tidak dapat dilebih-lebihkan. Seluruh mikroprosesor modern, dari prosesor sederhana pada kalkulator hingga GPU canggih pada superkomputer, beroperasi dengan memanipulasi sinyal listrik yang merepresentasikan nilai 0 dan 1. Setiap instruksi program, setiap perhitungan aritmetika, dan setiap keputusan logis yang dilakukan komputer pada akhirnya direduksi menjadi serangkaian operasi aljabar Boolean yang dieksekusi oleh jutaan transistor pada chip silikon \cite{allaboutcircuits_boolean}. Oleh karena itu, penguasaan aljabar Boolean merupakan kompetensi fundamental yang harus dimiliki oleh setiap profesional di bidang teknik informatika dan ilmu komputer.
